\section{Introduzione}
La simulazione del movimento collettivo dei pedoni rappresenta uno dei casi più interessanti di intersezione tra sistemi multi-agente, fisica statistica e progettazione urbana. In particolare, i nodi di intersezione ad alta densità costituiscono configurazioni critiche in cui il sistema può transitare rapidamente da un regime di \textit{laminar flow} a uno stato congestionato, con perdita di throughput e incremento del tempo di permanenza \cite{zhang2012}.

Nel presente lavoro si considera una geometria a rotonda pedonale a quattro bracci, implementata in NetLogo nel modello \texttt{DAI\_roundabout.nlogo}. La tesi principale è che la rotonda possa essere interpretata come un dispositivo di \textit{Distributed Artificial Intelligence passiva}: il coordinamento non deriva da regole comunicative esplicite tra agenti, ma dall'architettura dello spazio e dai campi locali che mediano la scelta dell'azione successiva.

Questa interpretazione è coerente con la nozione di \textit{stigmergia ambientale}. L'ambiente incorpora informazione decisionale sotto forma di gradienti statici verso la destinazione, tracce dinamiche lasciate dal traffico recente e costi repulsivi legati alla prossimità. In tal modo, la geometria non è un vincolo esterno ma parte attiva del processo computazionale distribuito.

La letteratura di riferimento suggerisce quattro assi teorici principali. Burstedde et al. introducono il \textit{Floor Field Model} come formalismo discreto per descrivere il moto pedonale con esclusione di volume e attrazione verso campi statici e dinamici \cite{burstedde2001}. Helbing et al. mostrano come lane formation e auto-organizzazione possano emergere da interazioni locali \cite{helbing2005}. Weng et al. e Yanagisawa et al. evidenziano che la rotonda, se ben progettata, riduce i conflitti frontali tipici degli incroci lineari \cite{weng2006,yanagisawa2019}. Zhang fornisce invece il quadro sperimentale per interpretare la relazione tra densità, flusso e velocità \cite{zhang2012}.
